\subsection{Lokale Extrema}

% --- Taylor-Formel in höheren Dimensionen ---
\begin{satz}{}
(Taylor-Formel)  
Es sei \(D \subset \mathbb{R}^n\) offen und \( f \in C^m(D) \). Sei \( x_0 \in D \) mit \( [x_0, x] \subset D \). Dann existiert \( \xi \in [x_0, x] \) mit


\[
f(x) = f(x_0) + \sum_{i=1}^{n} \frac{\partial f}{\partial x_i}(x_0) (x_i - x_{0,i}) + \frac{1}{2} \sum_{i,j=1}^{n} \frac{\partial^2 f}{\partial x_i \partial x_j} (\xi) (x_i - x_{0,i}) (x_j - x_{0,j}).
\]


\end{satz}

\begin{beweisbox}
Die Aussage folgt durch mehrdimensionale Taylorentwicklung mit Restglied. Der Beweis orientiert sich an der Einhüllenden der Funktion und nutzt die Darstellungen mittels Differentialoperatoren.  
% Details zur vollständigen Herleitung sind nicht vollständig erkennbar.
\end{beweisbox}

\begin{definition}{}
Eine Funktion \( f: D \to \mathbb{R} \), wobei \( D \subset \mathbb{R}^n \) offen ist, hat in \( x_0 \in D \) ein **lokales Maximum**, falls eine offene Umgebung \( U \) von \( x_0 \) existiert, sodass


\[
\forall x \in U: \quad f(x) \leq f(x_0).
\]


Ein **lokales Minimum** wird analog definiert.
\end{definition}

\begin{lemma}
Ist \( f: D \to \mathbb{R} \) partiell differenzierbar und hat in \( x_0 \) ein lokales Extremum (Minimum oder Maximum), so gilt:


\[
\operatorname{grad} f(x_0) = 0.
\]


Das bedeutet:


\[
\frac{\partial f}{\partial x_i} (x_0) = 0 \quad \forall i=1,\dots,n.
\]


\end{lemma}

\begin{beweisbox}
Falls \( x_0 \) ein lokales Extremum ist, existiert eine Umgebung, sodass \( f(x) \) dort entweder nie größer oder nie kleiner als \( f(x_0) \) ist. Dies führt durch das Richtungsableitungsprinzip zur Notwendigkeit, dass alle ersten Ableitungen verschwinden.
\end{beweisbox}

% --- Satz von Schwarz ---
\begin{satz}{}
(Satz von Schwarz)  
Sei \( f \in C^2(D) \), dann ist die Hesse-Matrix


\[
H_f(x) = \Bigg( \frac{\partial^2 f}{\partial x_i \partial x_j}(x) \Bigg)_{i,j}
\]


symmetrisch, d.h. es gilt für alle \( i,j \)


\[
\frac{\partial^2 f}{\partial x_i \partial x_j} = \frac{\partial^2 f}{\partial x_j \partial x_i}.
\]


\end{satz}

\begin{definition}{}
Eine quadratische Form \( Q(x) = x^T A x \) ist:
- **positiv definit**, falls \( Q(x) > 0 \) für alle \( x \neq 0 \).
- **negativ definit**, falls \( Q(x) < 0 \) für alle \( x \neq 0 \).
- **indefinit**, falls es \( x \neq 0 \) gibt mit \( Q(x) > 0 \) und \( Q(y) < 0 \) für manche \( y \neq 0 \).
\end{definition}

\begin{satz}{}
Sei \( f \in C^2(D) \) mit \( \operatorname{grad} f(x_0) = 0 \). Dann gilt:
\begin{enumerate}
    \item Falls die Hesse-Matrix \( H_f(x_0) \) **positiv definit** ist, besitzt \( f \) in \( x_0 \) ein lokales Minimum.
    \item Falls \( H_f(x_0) \) **negativ definit** ist, besitzt \( f \) in \( x_0 \) ein lokales Maximum.
    \item Falls \( H_f(x_0) \) **indefinit** ist, existiert kein lokales Extremum.
\end{enumerate}
\end{satz}

\begin{beweisbox}
Die Aussage folgt durch Taylor-Entwicklung und die Analyse der quadratischen Form, welche durch die Hesse-Matrix definiert wird.  
% Einige Abschätzungen im Originaltext konnten nicht vollständig rekonstruiert werden.
\end{beweisbox}

\begin{beispielbox}
\begin{enumerate}
    \item Sei \( f(x,y) = x^2 + y^2 \). Dann ist
    

\[
    H_f(x,y) = \begin{pmatrix} 2 & 0 \\ 0 & 2 \end{pmatrix}.
    \]


    Da die Hesse-Matrix positiv definit ist, besitzt \( f \) in \( (0,0) \) ein lokales Minimum.
    
    \item Sei \( f(x,y) = -x^2 - y^2 \). Dann ist
    

\[
    H_f(x,y) = \begin{pmatrix} -2 & 0 \\ 0 & -2 \end{pmatrix}.
    \]


    Da die Hesse-Matrix negativ definit ist, besitzt \( f \) in \( (0,0) \) ein lokales Maximum.
    
    \item Sei \( f(x,y) = x^2 - y^2 \). Dann ist
    

\[
    H_f(x,y) = \begin{pmatrix} 2 & 0 \\ 0 & -2 \end{pmatrix}.
    \]


    Da die Hesse-Matrix indefinit ist, existiert in \( (0,0) \) kein lokales Extremum.
\end{enumerate}
\end{beispielbox}
